\section{Cálculo de costos}
\label{sec:calculo_costos}

\subsection{Costos de inversión inicial}

Comprende los activos tecnológicos, herramientas de prueba y materiales de campo indispensables para la formulación, diseño y puesta en marcha del software. En la Tabla~\ref{tab:costos_inversion} se desglosa el presupuesto de inversión consolidado.

\begin{table}[htbp]
    \centering
    \caption{Costos de inversión del proyecto.}
    \label{tab:costos_inversion}
    \renewcommand{\arraystretch}{1.35}
    \begin{tabularx}{\textwidth}{L C C C C}
        \toprule
        \textbf{Concepto} & \textbf{Cantidad} & \textbf{Unidad} & \textbf{Costo unit. (Bs)} & \textbf{Costo total (Bs)} \\
        \midrule
        Computadoras portátiles & 2 & unidad & 4300.00 & 8600.00 \\
        Teléfonos celulares & 2 & unidad & 1200.00 & 2400.00 \\
        Impresora multifuncional & 1 & unidad & 2100.00 & 2100.00 \\
        Internet (2 mensualidades) & 2 & mensualidad & 140.00 & 280.00 \\
        Hojas de papel bond & 100 & hoja & 1.00 & 100.00 \\
        Transporte local & 9 & pasaje & 5.00 & 45.00 \\
        Cuadernos & 2 & unidad & 18.00 & 36.00 \\
        Archivadores & 2 & unidad & 4.00 & 8.00 \\
        Lapiceros & 5 & unidad & 1.00 & 5.00 \\
        \midrule
        \textbf{Total inversión} & & & & \textbf{13574.00} \\
        \bottomrule
    \end{tabularx}
    \notatabla{Elaboración propia. Total 13574.00 Bs con recursos propios, sin préstamos.}
\end{table}


\subsection{Costos operativos fijos y variables}

Los costos fijos mensuales ascienden a 190 Bs, correspondientes al servicio de internet de banda ancha (140 Bs) y fondo de mantenimiento de hardware (50 Bs). Los costos variables se estiman en 100 Bs por ciclo de inscripción para transporte local y papelería de respaldo administrativo.

\subsection{Análisis de impacto y sostenibilidad económica}

Al tratarse de una innovación de software público sin fines comerciales para el MÓDULO TECNOLÓGICO PRODUCTIVO SAN JULIÁN, la rentabilidad se mide en términos de retorno social y eficiencia institucional. El costo marginal por postulante registrado se reduce a menos de 0,50 Bs, mientras que para las familias de San Julián representa un ahorro conjunto superior a 20.000 Bs por año en pasajes y trámites presenciales, consolidando una solución plenamente sostenible.
